Abbreviate \(D=D_{J_y,J_z}\), \(M=2^{dJ_z}\), \(K=K_{J_y,J_z,n}\), \(\nu_k=\nu_{k,J_y,J_z,n}\), \(A=A_{J_y,J_z,n}\), and \(V=V_{L,n}\), and put
\[
b_n=(nV)^{1/2},\qquad S=S_n=A+b_n.                                         \tag{1}
\]
All constants below are uniform over the declared model class.

The uniform expansion gives
\[
\widehat\theta^{\mathrm{an}}_{J_y,J_z,n}-\theta_n
=\mathbb L_{J_y,J_z,n}+\mathbb U_{J_y,J_z,n}
+B^y_{J_y,n}+B^z_{J_z,n}+R_n,                                             \tag{2}
\]
where
\[
R_n=o_P\!\left(\frac{\rho_n}{\sqrt n}+\frac A n
+2^{-2(s_y+1)J_y}+2^{-2s_zJ_z}\right).                                   \tag{3}
\]
Since \(V=n\operatorname{Var}(\mathbb L_{J_y,J_z,n})\), that theorem also gives
\[
b_n\asymp\rho_n\sqrt n.                                                   \tag{4}
\]
The assumed bias inequality and (1)--(4) imply \(nR_n=o_P(S)\).  The deterministic biases are explicitly subtracted in the normalized statistic.

We compare the exact first projection with the linear part of the equality Hessian.  Let \(F_{a,n}(\cdot\mid z)\) denote the conditional distribution functions.  Positivity makes each inverse distribution function \(c^{-1}\)-Lipschitz, and
\[
\|F_{1,n}(\cdot\mid z)-F_{0,n}(\cdot\mid z)\|_\infty
\le\|\Delta_n(\cdot,z)\|_\infty.                                        \tag{5}
\]
Expanding the one-dimensional inverse maps in (5), using the zero outcome mass, yields the Neumann linearization
\[
-\partial_y\{\bar q_n(y,z)\partial_yu_n(y,z)\}=\Delta_n(y,z),
\qquad \partial_yu_n(0,z)=\partial_yu_n(1,z)=0,                            \tag{6}
\]
so \(u_n=\mathcal G_{n,z}\Delta_n(\cdot,z)\) after imposing zero integral.  A second expansion of the inverse identities is uniform because \(c\le q_{a,n}\le C\).  Centering the two arm derivatives and applying the wavelet tail estimates from the spectrum proposition gives
\[
\max_a\|s_{a,n}-s^{(2)}_{a,J_y,J_z,n}\|_{L^2(P_{a,n})}
\le C\rho_n\{\|\Delta_n\|_\infty
+2^{-(s_y+1)J_y}+2^{-s_zJ_z}\}.                                          \tag{7}
\]
Indeed, the nonlinear displacement is \(O(\|\Delta_n\|_\infty)\) by (5), and its product with the inverse-elliptic energy norm \(\rho_n\) bounds the first term.  The extra inverse-elliptic factor \(2^{-2j_y}\) and the anisotropic coefficient tails give the other two terms.

For fold \(e\), define
\[
\bar\Phi_{a,e}=m^{-1}\sum_{i\in I_{a,e}}\Phi_{J_y,J_z}(W_{a,i,n}),
\qquad
Z_{e,n}=\Gamma_{J_y,J_z,n}^{-1/2}\sqrt m
\{\bar\Phi_{1,e}-\bar\Phi_{0,e}-\delta_{J_y,J_z,n}\}.                    \tag{8}
\]
The two-arm sampling and covariance definition imply
\[
EZ_{e,n}=0,\qquad E(Z_{e,n}Z_{e,n}^\top)=I,
\qquad Z_{1,n}\ \text{and}\ Z_{2,n}\ \text{are independent}.            \tag{9}
\]
Put
\[
G_n=\sum_{e=1}^2 2\sqrt m\,\delta_{J_y,J_z,n}^\top
H_{J_y,J_z,n}\Gamma_{J_y,J_z,n}^{1/2}Z_{e,n}.                             \tag{10}
\]
Equations (4), (7), and (9), together with \(\|\Delta_n\|_\infty\to0\) and the vanishing projection tails, give
\[
\frac{\|n\mathbb L_{J_y,J_z,n}-G_n\|_{L^2}}S=o(1).                       \tag{11}
\]
If \(\rho_n=0\), positivity of the inverse form implies \(\Delta_n=0\), so both numerators in (11) are zero.  Diagonalizing \(K\) and using the definition of \(\eta_{k,n}\) gives
\[
\operatorname{Var}(G_n)=8\sum_k\nu_k^2\eta_{k,n}^2=nV+o(S^2).            \tag{12}
\]

We next prove universality of the canonical quadratic projection.  The upper inverse-form bound and compact support give \(\rho_n\le C\|\Delta_n\|_\infty\), hence \(\rho_n\to0\).  The spectrum proposition gives
\[
A\asymp M^{1/2},\qquad M\le D=o(n^2).                                    \tag{13}
\]
Equations (1), (4), and (13) imply \(S/n\to0\).  Since the premise includes \(2^{-2s_zJ_z}=O(S/n)\), necessarily \(J_z\to\infty\), and therefore
\[
M\to\infty,
\qquad
\frac{\|K\|_{\mathrm{op}}}{\|K\|_{\mathrm{HS}}}\lesssim M^{-1/2}\to0. \tag{14}
\]

Index all evaluation observations by a single index \(r\), retaining fold and arm labels.  Diagonal removal and canonical centering give
\[
Q_n:=n\mathbb U_{J_y,J_z,n}=\sum_{r<s}h_{n,rs}(W_r,W_s),
\qquad E\{h_{n,rs}\mid W_r\}=E\{h_{n,rs}\mid W_s\}=0.                  \tag{15}
\]
The Green kernel of (6) is uniformly bounded in its outcome arguments.  Compact support of the covariate wavelets gives a projection envelope \(\kappa_{J_z}\) with
\[
\iint\kappa_{J_z}(z,z')^2dR_n(z)dR_n(z')\le CM,
\qquad
\iint\kappa_{J_z}(z,z')^4dR_n(z)dR_n(z')\le CM^3.                        \tag{16}
\]
At each scale only finitely many wavelets overlap; summing their squared magnitudes gives the first bound, and repeating the overlap count after squaring gives the second.  Covariance whitening changes these estimates only by fixed factors.  Pairing indices in (15) yields
\[
\begin{aligned}
&\operatorname{Var}(Q_n)=A^2\{1+o(1)\},
\qquad \max_r\sum_{s\ne r}Eh_{n,rs}^2\le C M/m,\\
&\sum_{r<s}Eh_{n,rs}^4\le CM^3/m^2,
\qquad
\frac{\operatorname{tr}(K^4)}{\{\operatorname{tr}(K^2)\}^2}
\le\frac{\|K\|_{\mathrm{op}}^2}{\|K\|_{\mathrm{HS}}^2}\le C/M.
\end{aligned}                                                             \tag{17}
\]
Thus the maximum row variance divided by \(A^2\) is \(O(m^{-1})\), and the fourth-moment sum divided by \(A^4\) is \(O(M/m^2)=o(1)\), using \(M\le D=o(n^2)\).

Order the independent observations, let \(\mathcal F_r=\sigma(W_1,\ldots,W_r)\), write \(G_n=\sum_r\ell_{n,r}(W_r)\), and define
\[
\mathcal D_{n,r}=\ell_{n,r}(W_r)+\sum_{s<r}h_{n,sr}(W_s,W_r).              \tag{18}
\]
Canonicality makes this a martingale-difference array with sum \(G_n+Q_n\).  The inverse-map bound leading to (7) gives
\[
\omega_n:=\max_r\|\ell_{n,r}\|_\infty/S=o(1).                           \tag{19}
\]
Expanding conditional squares and fourth powers in (18), a nonzero expectation must pair every observation index at least twice.  Shared-index pairs give the row term in (17), repeated edges give its fourth-moment term, four-edge cycles give \(\operatorname{tr}(K^4)\), and terms with a linear factor are bounded by (19).  Consequently, for every \(\epsilon>0\),
\[
\sum_rE\left[\left(\frac{\mathcal D_{n,r}}S\right)^2
\mathbf1\{|\mathcal D_{n,r}|>\epsilon S\}\right]\to0,                  \tag{20}
\]
and
\[
\frac{E\left[\left\{\sum_rE(\mathcal D_{n,r}^2\mid\mathcal F_{r-1})
-E(G_n+Q_n)^2\right\}^2\right]}{S^4}
\le C\left\{m^{-1}+M/m^2+M^{-1}+\omega_n^2\right\}\to0.                \tag{21}
\]
Conditional Taylor expansion of \(e^{itx}\), with (20) controlling its summed remainder and (21) replacing the conditional variance by its expectation, now gives
\[
E\exp\{it(G_n+Q_n)/S\}
-\exp\left[-\frac{t^2}{2S^2}\{A^2+nV+o(S^2)\}\right]\to0.              \tag{22}
\]
The covariance of \(G_n\) and \(Q_n\) is zero by (15), and (12) supplies the linear variance.

On the declared Gaussian array define
\[
\mathcal Q_n^G=\sum_{e=1}^2\sum_k\nu_k
\{(\xi_{e,k}+\eta_{k,n})^2-(1+\eta_{k,n}^2)\}.                           \tag{23}
\]
Put \(c_{e,k}=\nu_k/S\) and \(l_{e,k}=2\nu_k\eta_{k,n}/S\).  Direct one-dimensional Gaussian integration and independence give
\[
\log E e^{it\mathcal Q_n^G/S}
=\sum_{e,k}\left[-itc_{e,k}-\frac12\log(1-2itc_{e,k})
-\frac{t^2l_{e,k}^2}{2(1-2itc_{e,k})}\right].                             \tag{24}
\]
By (12), (14), and \(\sum_{e,k}l_{e,k}^2=\{nV+o(S^2)\}/S^2\), the Taylor remainder in (24) is at most
\[
C_t\max_k\frac{|\nu_k|}{S}
\left\{\frac{A^2}{S^2}+\sum_{e,k}l_{e,k}^2\right\}=o(1).                \tag{25}
\]
Thus (24) equals the exponent in (22), up to \(o(1)\).  Every subsequence contains a further subsequence on which \(A/S\) converges; on it, (22) and (24) converge to the same centered normal distribution, whose variance lies in \([1/2,1]\).  This subsequence argument proves
\[
d_{\mathrm{BL}}\{\mathcal L((G_n+Q_n)/S),
\mathcal L(\mathcal Q_n^G/S)\}\to0.                                     \tag{26}
\]
Combining (2), \(nR_n=o_P(S)\), (11), and (26) proves the stated triangular approximation.

At equality, positivity of the inverse form gives \(\Delta_n=0\); hence \(\delta=0\), every \(\eta_{k,n}=0\), \(V=0\), and \(\mathbb L=0\).  Equations (14), (24), and (25) yield
\[
\frac1A\sum_{e,k}\nu_k(\xi_{e,k}^2-1)\Rightarrow\mathcal N(0,1),
\]
and (26) transfers the conclusion to \(n\mathbb U/A\).

For the qualified chaos assertion, put \(w_{k,n}=\nu_k/A\), pad the arrays by zeros, and suppose \(w_n\to w\) and \(\eta_n\to\eta\) in \(\ell^2\) along a subsequence.  Coupling on one Gaussian array gives \(L^2\) convergence because the centered quadratic difference has variance \(4\|w_n-w\|_2^2\), while
\[
\|w_n\eta_n-w\eta\|_2
\le\|w_n-w\|_2\|\eta_n\|_\infty
+\|w\|_\infty\|\eta_n-\eta\|_2\to0.                                  \tag{27}
\]
The limiting series is therefore a well-defined centered Gaussian chaos, and (26) transfers it to the estimator after taking a further subsequence if \(S/A\) requires a limit.  Conversely, in the diffuse null regime, \(\max_k|w_{k,n}|\to0\) but \(\|w_n\|_2=1/2\), so no strongly convergent \(\ell^2\) subsequence exists; this proves the stated qualification.

Under \(D\log n/m=o(1)\), the pilot-control lemma gives covariance and Hessian consistency.  Its coercive resolvent estimate is relative in the elliptic energy, so for \(\bar K_n^{\mathrm{obs}}=\{\widehat K^{(-1)}+\widehat K^{(-2)}\}/2\),
\[
|x^\top(\bar K_n^{\mathrm{obs}}-K)y|
\le o_P(1)(x^\top Kx)^{1/2}(y^\top Ky)^{1/2}.
\]
Summing this inequality in an eigenbasis gives
\[
\|\bar K_n^{\mathrm{obs}}-K\|_{\mathrm{HS}}
=o_P(\|K\|_{\mathrm{HS}})=o_P(A).                                        \tag{28}
\]
For a common standard Gaussian vector \(Z\), direct Gaussian moment expansion gives the conditional coupling bound
\[
E\left[\{Z^\top(\bar K_n^{\mathrm{obs}}-K)Z
-\operatorname{tr}(\bar K_n^{\mathrm{obs}}-K)\}^2\mid\text{pilots}\right]
=2\|\bar K_n^{\mathrm{obs}}-K\|_{\mathrm{HS}}^2=o_P(S^2),
\]
which proves the observed-kernel replacement.

Finally, at a fixed unequal law, the exact centered first-order score is bounded on the compact support because the cost is bounded and the densities are uniformly positive.  Expanding its characteristic function exactly as in (20)--(22), now without the quadratic summand, gives
\[
\mathbb L_{J_y,J_z,n}/\sqrt{V/n}\Rightarrow\mathcal N(0,1).
\]
The hypothesis \(\sqrt{nV}/A\to\infty\) makes \(\mathbb U=o_P(\sqrt{V/n})\); the stated bias and remainder condition removes all other terms in (2).  Since \(\widehat V_{L,n}/V_{L,n}\to_P1\), multiplication by \(\sqrt{V_{L,n}/\widehat V_{L,n}}\) proves the final studentized normal limit.